The Gaussian calculation uses the Hermitian amplitude/imaginary-pairing matrix in Eqs.~\eqref{eq:kernelAA}--\eqref{eq:kernelAP}. A second representation linearizes the spatial quasiclassical equations through two Riccati coherence sectors~\cite{Eilenberger1968,Eschrig2000}. Uniform-amplitude, phase Ward, twist--amplitude, current, and zero-field controls establish a common normalization before the spatial comparison.

For the nonlinear response, both causal coherence sectors are expanded in the signed perturbation amplitude, retaining every generated Fourier convolution mode. Spatial averaging takes the zero Fourier coefficient of the retained polynomial; omitted orders and modes remain in the remainder estimates. The normalized anomalous force is integrated along the pairing path and added to the second linear coefficient, avoiding subtraction of large independently accumulated thermodynamic potentials. Sign pairing removes odd orders before absolute errors are formed. Appendix~\ref{app:response} gives the recurrence, both normalization denominators, and the path formula.

The uncertainty construction distinguishes endpoint uncertainty, numerical resolution, and finite-amplitude corrections. The empirical coordinate rectangle is propagated through the background, both coherence sectors, the force, and the temperature prefactor. Low frequencies retain profile-residual and inverse estimates over a connected branch neighbourhood; high frequencies use an analytic complex-amplitude disk and an infinite-frequency tail envelope covering all higher even orders. Neither a finite cutoff nor the retained Fourier polynomial is treated as the full response.

Complete numerical allowances add the linear baseline error, primitive arithmetic estimates, profile and force residuals, path discrepancy, frequency tail, and single- and mixed-refinement contributions. The endpoint allowance is added once. These deterministic additive allowances define the empirical upper endpoints; they are neither statistical confidence intervals nor rigorous interval enclosures. The sign criteria require agreement between representations, negative upper endpoints, and the predefined refinement and source checks. The quadratic total allowances must also be at most one quarter of the coefficient magnitudes.

The angular-origin comparison shares the endpoint, second linear baseline, and continuous-domain estimates. It probes angular discretization sensitivity rather than supplying independent physical observations. Exact refinement levels, source allocations, arithmetic tolerances, and complete acceptance accounting appear in Appendices~\ref{app:response}--\ref{app:repro}.
